% Emacs, this is -*-latex-*-

\ex{I-equival\^encia}
\label{ex:I-equiv}

\begin{exenumerate}
    \item Qual dos tr\^es grafos representa o mesmo conjunto de independ\^encias? Explique.\\
    
    \begin{tikzpicture}[dgraph]
        \node[cont] (v) at (0,0) {$v$};
        \node[cont, below right= of v] (w) {$w$};
        \node[cont, right= of v] (x) {$x$};
        \node[cont, below right= of x] (y) {$y$};
        \node[cont, right= of x] (z) {$z$};
        
        \node[below= 0.2 of w] {Grafo 1};
        
        \draw (v) -- (w);
        \draw (w) -- (x);
        \draw (x) -- (y);
        \draw (z) -- (y);   
    \end{tikzpicture}
    \hspace{5ex}
    \begin{tikzpicture}[dgraph]
        \node[cont] (v) at (0,0) {$v$};
        \node[cont, below right= of v] (w) {$w$};
        \node[cont, right= of v] (x) {$x$};
        \node[cont, below right= of x] (y) {$y$};
        \node[cont, right= of x] (z) {$z$};
        
        \node[below= 0.2 of w] {Grafo 2};
        
        \draw (w) -- (v);
        \draw (x) -- (w);
        \draw (x) -- (y);
        \draw (z) -- (y);   
    \end{tikzpicture}
    \hspace{5ex}
    \begin{tikzpicture}[dgraph]
        \node[cont] (v) at (0,0) {$v$};
        \node[cont, below right= of v] (w) {$w$};
        \node[cont, right= of v] (x) {$x$};
        \node[cont, below right= of x] (y) {$y$};
        \node[cont, right= of x] (z) {$z$};
        
        \node[below= 0.2 of w] {Grafo 3};
        
        \draw (w) -- (v);
        \draw (x) -- (w);
        \draw (y) -- (x);
        \draw (z) -- (y);
        
    \end{tikzpicture}
    
    \begin{solution}
        Para verificar se os grafos s\~ao I-equivalentes, temos que verificar os esqueletos e as imoralidades. Todos possuem o mesmo esqueleto, mas o grafo 1 e o grafo 2 tamb\'em possuem a mesma imoralidade. A resposta \'e, portanto: os grafos 1 e 2 codificam as mesmas independ\^encias.\\
        \begin{center}
            \begin{tikzpicture}[ugraph]
                \node[cont] (v) at (0,0) {$v$};
                \node[cont, below right= of v] (w) {$w$};
                \node[cont, right= of v] (x) {$x$};
                \node[cont, below right= of x] (y) {$y$};
                \node[cont, right= of x] (z) {$z$};
                
                \node[below= 0.2 of w] {esqueleto};
                
                \draw (w) -- (v);
                \draw (x) -- (w);
                \draw (y) -- (x);
                \draw (z) -- (y);
                
            \end{tikzpicture}
            \hspace{4ex}
            \begin{tikzpicture}[ugraph]
                \node[cont] (v) at (0,0) {$v$};
                \node[cont, below right= of v] (w) {$w$};
                \node[cont, right= of v] (x) {$x$};
                \node[cont, below right= of x] (y) {$y$};
                \node[cont, right= of x] (z) {$z$};
                
                \node[below= 0.2 of w] {imoralidade};
                
                \draw (w) -- (v);
                \draw (x) -- (w);
                \draw[->, red] (x) -- (y);
                \draw[->,red] (z) -- (y);
                
            \end{tikzpicture}
        \end{center}
    \end{solution}
    
    
    \item Qual dos tr\^es grafos representa o mesmo conjunto de independ\^encias? Explique.\\
    
    \begin{tikzpicture}[dgraph]
        \node[cont] (v) at (0,0) {$v$};
        \node[cont, right= of v] (x) {$x$};
        \node[cont, below = of x] (w) {$w$};
        \node[cont, right= of w] (y) {$y$};
        \node[cont, below= of y] (z) {$z$};
        
        \node[below= 2 of w] {Grafo 1};
        
        \draw (v) -- (x);
        \draw (v) -- (w);
        \draw (x) -- (w);
        \draw (w) -- (z);
        \draw (y) -- (z);   
    \end{tikzpicture}
    \hspace{5ex}
    \begin{tikzpicture}[dgraph]
        \node[cont] (v) at (0,0) {$v$};
        \node[cont, right= of v] (x) {$x$};
        \node[cont, below = of x] (w) {$w$};
        \node[cont, right= of w] (y) {$y$};
        \node[cont, below= of y] (z) {$z$};
        
        \node[below= 2 of w] {Grafo 2};
        
        \draw (x) -- (v);
        \draw (w) -- (v);
        \draw (x) -- (w);
        \draw (w) -- (z);
        \draw (y) -- (z);   
    \end{tikzpicture}
    \hspace{5ex}
    \begin{tikzpicture}[dgraph]
        \node[cont] (v) at (0,0) {$v$};
        \node[cont, right= of v] (x) {$x$};
        \node[cont, below = of x] (w) {$w$};
        \node[cont, right= of w] (y) {$y$};
        \node[cont, below= of y] (z) {$z$};
        
        \node[below= 2 of w] {Grafo 3};
        
        \draw (v) -- (w);
        \draw (x) -- (w);
        \draw (w) -- (z);
        \draw (y) -- (z);   
    \end{tikzpicture}
    \hspace{5ex}
    
    \begin{solution}
        O esqueleto do grafo 3 \'e diferente do esqueleto dos grafos 1
        e 2, portanto o grafo 3 n\~ao pode ser I-equivalente aos grafos 1 ou 2, e
        n\~ao precisamos verificar as imoralidades para o grafo 3. Os grafos 1 e 2
        possuem o mesmo esqueleto e tamb\'em possuem a mesma
        imoralidade. Logo, os grafos 1 e 2 s\~ao I-equivalentes. Note que o n\'o
        $w$ no grafo 1 est\'a em uma configura\'c\~ao de colisor ao longo da trilha $v-w-x$,
        mas n\~ao \'e uma imoralidade porque seus pais est\~ao conectados
        (aresta de cobertura); de forma equivalente para o n\'o $v$ no grafo 2.
        
        \begin{center}
            \begin{tikzpicture}[ugraph]
                \node[cont] (v) at (0,0) {$v$};
                \node[cont, right= of v] (x) {$x$};
                \node[cont, below = of x] (w) {$w$};
                \node[cont, right= of w] (y) {$y$};
                \node[cont, below= of y] (z) {$z$};
                
                \node[below= 2 of w] {esqueleto};
                
                \draw (v) -- (x);
                \draw (v) -- (w);
                \draw (x) -- (w);
                \draw (w) -- (z);
                \draw (y) -- (z);   
            \end{tikzpicture}
            \hspace{4ex}
            \begin{tikzpicture}[ugraph]
                \node[cont] (v) at (0,0) {$v$};
                \node[cont, right= of v] (x) {$x$};
                \node[cont, below = of x] (w) {$w$};
                \node[cont, right= of w] (y) {$y$};
                \node[cont, below= of y] (z) {$z$};
                
                \node[below= 2 of w] {imoralidade};
                
                \draw (v) -- (x);
                \draw (v) -- (w);
                \draw (x) -- (w);
                \draw[->, red] (w) -- (z);
                \draw[->, red] (y) -- (z);   
            \end{tikzpicture}
        \end{center}
    \end{solution}
    
    \item Assuma que o grafo abaixo seja um mapa perfeito para um conjunto de independ\^encias $\mathcal{U}$.
    
    \begin{center}
        \begin{tikzpicture}[dgraph]
            \node[cont] (x1) at (0,0) {$x_1$};
            \node[cont, right= of x1] (x2) {$x_2$};
            \node[cont, below = of x1] (x3) {$x_3$};
            \node[cont, right= of x3] (x4) {$x_4$};
            \node[cont, right= of x4] (x5) {$x_5$};
            \node[cont, below= of x3] (x6) {$x_6$};
            \node[cont, below= of x4] (x7) {$x_7$};
            
            \node[below= 0.2 of x6] {Grafo 0};
            
            \draw (x1) -- (x3);
            \draw (x2) -- (x4);
            \draw (x2) -- (x5);
            \draw (x3) -- (x6);
            \draw (x4) -- (x7);
            \draw (x5) -- (x7);
            \draw (x7) -- (x6);
            \draw (x3) -- (x7);
        \end{tikzpicture}
    \end{center}
    
    Para cada um dos tr\^es grafos abaixo, explique se o grafo \'e um mapa perfeito, um I-mapa, ou n\~ao \'e um I-mapa para $\mathcal{U}$.
    
    \begin{center}
        \begin{tikzpicture}[dgraph]
            \node[cont] (x1) at (0,0) {$x_1$};
            \node[cont, right= of x1] (x2) {$x_2$};
            \node[cont, below = of x1] (x3) {$x_3$};
            \node[cont, right= of x3] (x4) {$x_4$};
            \node[cont, right= of x4] (x5) {$x_5$};
            \node[cont, below= of x3] (x6) {$x_6$};
            \node[cont, below= of x4] (x7) {$x_7$};
            
            \node[below= 0.2 of x6] {Grafo 1};
            
            \draw (x1) -- (x3);
            \draw (x2) -- (x4);
            \draw (x2) -- (x5);
            \draw (x3) -- (x6);
            \draw (x4) -- (x7);
            \draw (x7) -- (x5);
            \draw (x7) -- (x6);
            \draw (x3) -- (x7);
        \end{tikzpicture}
        \begin{tikzpicture}[dgraph]
            \node[cont] (x1) at (0,0) {$x_1$};
            \node[cont, right= of x1] (x2) {$x_2$};
            \node[cont, below = of x1] (x3) {$x_3$};
            \node[cont, right= of x3] (x4) {$x_4$};
            \node[cont, right= of x4] (x5) {$x_5$};
            \node[cont, below= of x3] (x6) {$x_6$};
            \node[cont, below= of x4] (x7) {$x_7$};
            
            \node[below= 0.2 of x6] {Grafo 2};
            
            \draw (x1) -- (x3);
            \draw (x2) -- (x4);
            \draw (x2) -- (x5);
            \draw (x3) -- (x6);
            \draw (x4) -- (x7);
            \draw (x5) -- (x7);
            \draw (x7) -- (x6);
            \draw (x7) -- (x3);
        \end{tikzpicture}
        \begin{tikzpicture}[dgraph]
            \node[cont] (x1) at (0,0) {$x_1$};
            \node[cont, right= of x1] (x2) {$x_2$};
            \node[cont, below = of x1] (x3) {$x_3$};
            \node[cont, right= of x3] (x4) {$x_4$};
            \node[cont, right= of x4] (x5) {$x_5$};
            \node[cont, below= of x3] (x6) {$x_6$};
            \node[cont, below= of x4] (x7) {$x_7$};
            
            \node[below= 0.2 of x6] {Grafo 3};
            
            \draw (x1) -- (x3);
            \draw (x2) -- (x4);
            \draw (x5) -- (x2);
            \draw (x3) -- (x6);
            \draw (x4) -- (x7);
            \draw (x5) -- (x7);
            \draw (x7) -- (x6);
            \draw (x3) -- (x7);
        \end{tikzpicture}
    \end{center}
    
    \begin{solution}
        \begin{itemize}
            \item O grafo 1 possui uma imoralidade $x_2 \rightarrow x_5 \leftarrow x_7$ que o grafo 0 n\~ao possui. O grafo n\~ao \'e, portanto, I-equivalente ao grafo 0 e n\~ao pode ser um mapa perfeito. Al\'em disso, o grafo 1 afirma que $x_2 \independent x_7 | x_4$, o que n\~ao \'e o caso para o grafo 0. Como o grafo 0 \'e um mapa perfeito para $\mathcal{U}$, o grafo 1 afirma uma independ\^encia que n\~ao vale para $\mathcal{U}$ e, portanto, n\~ao pode ser um I-mapa para $\mathcal{U}.$
            \item O grafo 2 possui uma imoralidade $x_1 \rightarrow x_3 \leftarrow x_7$ que o grafo 0 n\~ao possui. O grafo 2 afirma, assim, que $x_1 \independent x_7$, o que n\~ao \'e o caso para o grafo 0. Logo, pela mesma raz\~ao do grafo 1, o grafo 2 n\~ao \'e um I-mapa para $\mathcal{U}$.
            \item O grafo 3 possui o mesmo esqueleto e o mesmo conjunto de imoralidades que o grafo 0. Ele \'e, portanto, I-equivalente ao grafo 0 e, consequentemente, tamb\'em um mapa perfeito.
        \end{itemize}
    \end{solution}
    
    
\end{exenumerate}


\ex{I-mapas m\'inimos}
\label{ex:minimal-I-maps-1}
\begin{exenumerate}
    \item Assuma que o grafo $G$ na Figura \ref{fig:I-map} seja um I-mapa perfeito para $p(a,z,q,e,h)$. Determine o I-mapa direcionado m\'inimo usando a ordena\'c\~ao $(e,h,q,z,a)$. O grafo obtido \'e I-equivalente a $G$?
    \begin{figure}[ht]
        \begin{center}
            \scalebox{0.9}{ % x,y
                \begin{tikzpicture}[dgraph]
                    \node[cont] (a) at (0,2) {$a$};
                    \node[cont] (z) at (2,2) {$z$};
                    \node[cont] (q) at (1,1) {$q$};
                    \node[cont] (e) at (1,-0.4) {$e$};
                    \node[cont] (h) at (3,1) {$h$};
                    \draw(a) -- (q);
                    \draw(z) -- (h);
                    \draw(z) -- (q);
                    \draw(q) -- (e);
            \end{tikzpicture}}
        \end{center}
        \caption{\label{fig:I-map} I-mapa perfeito $G$ para o \exref{ex:minimal-I-maps-1}, quest\~ao \ref{q:directed-I-maps}.}
    \end{figure}
    \label{q:directed-I-maps}
    \begin{solution} Como o grafo $G$ \'e um I-mapa perfeito para $p$, podemos usar
        $G$ para verificar se $p$ satisfaz uma certa independ\^encia. Isso fornece a seguinte receita para construir o I-mapa direcionado m\'inimo:
        \begin{enumerate}
            \item Assuma uma ordena\'c\~ao das vari\'aveis. Denomine as vari\'aveis aleat\'orias ordenadas por $x_1, \ldots, x_d$.
            \item Para cada $i$, encontre um subconjunto m\'inimo de vari\'aveis $\pap_i \subseteq \pre_i$ tal que $$ x_i \independent \{\pre_i \setminus \pap_i \} \mid \pap_i$$ esteja em $\Ind(G)$ (funciona apenas se $G$ for um I-mapa perfeito para $\Ind(p)$)
            \item Construa um grafo com pais $\pa_i=\pap_i$.
        \end{enumerate}
        
        Nota: Para I-mapas $G$ que n\~ao s\~ao perfeitos, se o grafo n\~ao indicar que uma certa independ\^encia ocorre, temos que verificar se a independ\^encia de fato n\~ao ocorre para $p$. Se n\~ao o fizermos, n\~ao obteremos um I-mapa m\'inimo, mas apenas um I-mapa para $\Ind(p)$. Isso ocorre porque $p$ pode ter independ\^encias que n\~ao est\~ao codificadas no grafo $G$. 
        
        Dada a ordena\'c\~ao $(e,h,q,z,a)$, constru\'imos um grafo onde $e$ \'e a raiz. Pela Figura \ref{fig:I-map} (e pela suposi\'c\~ao de mapa perfeito), vemos que $h \independent e$ n\~ao ocorre. Definimos, assim, $e$ como pai de $h$, veja o primeiro grafo na Figura \ref{fig:I-map2}. Ent\~ao:
        \begin{itemize}
            \item Consideramos $q$: $\pre_q = \{e,h\}$. N\~ao h\'a subconjunto $\pap_q$ de $\pre_q$ no qual pud\'essemos condicionar para tornar $q$ independente de $\pre_q \setminus \pap_q$, de modo que definimos os pais de $q$ no grafo como $\pa_q = \{e,h\}$. (Segundo grafo na Figura \ref{fig:I-map2}.)
            \item Consideramos $z$: $\pre_z = \{e,h,q\}$. Pelo grafo na Figura \ref{fig:I-map}, vemos que para $\pap_z = \{q,h\}$ temos $z \independent \pre_z \setminus \pap_z | \pap_z$. Note que $\pap_z = \{q\}$ n\~ao funciona porque $z \independent e,h | q$ n\~ao ocorre. Definimos, assim, $\pa_z =\{q,h\}$. (Terceiro grafo na Figura \ref{fig:I-map2}.)
            \item Consideramos $a$: $\pre_a = \{e,h,q,z\}$. Este \'e o \'ultimo n\'o na ordena\'c\~ao. Para encontrar o conjunto m\'inimo $\pap_a$ para o qual $a \independent \pre_a \setminus \pap_a | \pap_a$, podemos determinar seu Markov blanket $\MB(a)$. O Markov blanket \'e o conjunto de pais (nenhum), filhos ($q$), e copais de $a$ ($z$) na Figura \ref{fig:I-map}, de modo que $\MB(a) = \{q,z\}$. Definimos, assim, $\pa_a = \{q,z\}$. (Quarto grafo na Figura \ref{fig:I-map2}.)
        \end{itemize}
        \begin{figure}[h]
            \centering
            \scalebox{0.9}{ % x,y
                \begin{tikzpicture}[dgraph]
                    \node[cont] (a) at (0,2) {$a$};
                    \node[cont] (z) at (2,2) {$z$};
                    \node[cont] (q) at (1,1) {$q$};
                    \node[cont] (e) at (1,-0.4) {$e$};
                    \node[cont] (h) at (3,1) {$h$};
                    \draw(e) -- (h);
                    
                    \draw[-,gray,dotted] (3.7,-0.4) -- (3.7,2);
            \end{tikzpicture}}
            \hspace{1ex}
            \scalebox{0.9}{ % x,y
                \begin{tikzpicture}[dgraph]
                    \node[cont] (a) at (0,2) {$a$};
                    \node[cont] (z) at (2,2) {$z$};
                    \node[cont] (q) at (1,1) {$q$};
                    \node[cont] (e) at (1,-0.4) {$e$};
                    \node[cont] (h) at (3,1) {$h$};
                    \draw(e) -- (h);
                    \draw(e) -- (q);
                    \draw(h) -- (q);
                    
                    \draw[-,gray,dotted] (3.7,-0.4) -- (3.7,2);
            \end{tikzpicture}}
            \hspace{1ex}
            \scalebox{0.9}{ % x,y
                \begin{tikzpicture}[dgraph]
                    \node[cont] (a) at (0,2) {$a$};
                    \node[cont] (z) at (2,2) {$z$};
                    \node[cont] (q) at (1,1) {$q$};
                    \node[cont] (e) at (1,-0.4) {$e$};
                    \node[cont] (h) at (3,1) {$h$};
                    \draw(e) -- (h);
                    \draw(e) -- (q);
                    \draw(h) -- (q);
                    \draw(q) -- (z);
                    \draw(h) -- (z);
                    
                    \draw[-,gray,dotted] (3.7,-0.4) -- (3.7,2);
            \end{tikzpicture}}
            \hspace{1ex}
            \scalebox{0.9}{ % x,y
                \begin{tikzpicture}[dgraph]
                    \node[cont] (a) at (0,2) {$a$};
                    \node[cont] (z) at (2,2) {$z$};
                    \node[cont] (q) at (1,1) {$q$};
                    \node[cont] (e) at (1,-0.4) {$e$};
                    \node[cont] (h) at (3,1) {$h$};
                    \draw(e) -- (h);
                    \draw(e) -- (q);
                    \draw(h) -- (q);
                    \draw(q) -- (z);
                    \draw(h) -- (z);
                    \draw(q) -- (a);
                    \draw(z) -- (a);    
            \end{tikzpicture}}
            \caption{\label{fig:I-map2} \exref{ex:minimal-I-maps-1}, Quest\~ao \ref{q:directed-I-maps}: Constru\'c\~ao de um I-mapa direcionado m\'inimo para a ordena\'c\~ao $(e,h,q,z,a)$. }
        \end{figure}
        
        Como o esqueleto no I-mapa m\'inimo obtido \'e diferente do esqueleto de $G$, n\~ao temos I-equival\^encia. Note que a ordena\'c\~ao $(e,h,q,z,a)$ resulta em um grafo mais denso (Figura \ref{fig:I-map2}) do que o grafo na Figura \ref{fig:I-map}. Embora seja um I-mapa m\'inimo, o grafo n\~ao mostra, por exemplo, que $a \independent z$. Al\'em disso, a interpreta\'c\~ao causal dos dois grafos \'e diferente.
        
    \end{solution}
    
    \item Para a cole\'c\~ao de vari\'aveis aleat\'orias $(a,z,h,q,e)$, s\~ao fornecidos os seguintes Markov blankets para cada vari\'avel:
    \begin{itemize}
        \item \MB(a) = \{q,z\}
        \item \MB(z) = \{a,q,h\}
        \item \MB(h) = \{z\}
        \item \MB(q) = \{a,z,e\}
        \item \MB(e) = \{q\}
    \end{itemize}
    \begin{exenumerate}
        \item Desenhe o I-mapa m\'inimo n\~ao direcionado que representa as independ\^encias.
        \item Indique uma distribui\'c\~ao de Gibbs que satisfa\'ca as rela\'c\~oes de independ\^encia especificadas pelos Markov blankets.
    \end{exenumerate}
    
    \begin{solution} Conectar cada vari\'avel a todas as vari\'aveis em seu Markov blanket fornece o I-mapa m\'inimo n\~ao direcionado desejado. Note que os Markov blankets n\~ao s\~ao mutuamente disjuntos. \begin{figure}[h] \begin{center} \scalebox{0.9}{ % x,y
                    \begin{tikzpicture}[ugraph] \node[cont] (a) at (0,2) {$a$}; \node[cont] (z) at (2,2) {$z$}; \node[cont] (q) at (1,1) {$q$}; \node[cont] (e) at (1,-0.4) {$e$}; \node[cont] (h) at (3,1) {$h$}; \draw(a) -- (q); \draw(a) -- (z); \draw[-,gray,dotted] (3.7,-0.4) -- (3.7,2);
                        
                        \node[below = of e] {Ap\'os $\MB(a)$};
                    \end{tikzpicture}
                    \vspace{3ex}
                    \begin{tikzpicture}[ugraph]
                        \node[cont] (a) at (0,2) {$a$};
                        \node[cont] (z) at (2,2) {$z$};
                        \node[cont] (q) at (1,1) {$q$};
                        \node[cont] (e) at (1,-0.4) {$e$};
                        \node[cont] (h) at (3,1) {$h$};
                        \draw(a) -- (q);
                        \draw(a) -- (z);
                        \draw(z) -- (h);
                        \draw(z) -- (q);
                        \draw[-,gray,dotted] (3.7,-0.4) -- (3.7,2);
                        \node[below = of e] {Ap\'os $\MB(z)$};
                    \end{tikzpicture}
                    \vspace{3ex}
                    \begin{tikzpicture}[ugraph]
                        \node[cont] (a) at (0,2) {$a$};
                        \node[cont] (z) at (2,2) {$z$};
                        \node[cont] (q) at (1,1) {$q$};
                        \node[cont] (e) at (1,-0.4) {$e$};
                        \node[cont] (h) at (3,1) {$h$};
                        \draw(a) -- (q);
                        \draw(a) -- (z);
                        \draw(z) -- (h);
                        \draw(z) -- (q);
                        \draw(q) -- (e);
                        \node[below = of e] {Ap\'os $\MB(q)$};
                \end{tikzpicture}}
            \end{center}
        \end{figure}
        
        Para distribui\'c\~oes positivas, o conjunto de distribui\'c\~oes que satisfazem a propriedade de Markov local relativa a um grafo (como dado pelos Markov blankets) \'e o mesmo que o conjunto de distribui\'c\~oes de Gibbs que se fatorizam de acordo com o grafo. Dado o I-mapa, podemos agora encontrar facilmente a distribui\'c\~ao de Gibbs
        $$p(a,z,h,q,e) = \frac{1}{Z} \phi_1(a,z,q) \phi_2(q,e) \phi_3(z,h),$$
        onde os $\phi_i$ devem assumir valores positivos em seu dom\'inio. Note que usamos o clique m\'aximo $(a,z,q)$.
        
    \end{solution}
\end{exenumerate}


\ex{I-equival\^encia entre grafos direcionados e n\~ao direcionados}

\begin{exenumerate}
    \item Verifique se os dois grafos seguintes s\~ao I-equivalentes listando e comparando as independ\^encias que cada grafo implica.
    
    \begin{center}
        \scalebox{1}{ % x,y
            \begin{tikzpicture}[ugraph]
                \node[cont] (z) at (0,0) {$z$};
                \node[cont] (y) at (1,1) {$y$};
                \node[cont] (x) at (0,2) {$x$};
                \node[cont] (u) at (-1,1) {$u$};
                \draw(x) -- (y);
                \draw(y) -- (z);
                \draw(z) -- (u);
                \draw(u) -- (x);
                \draw(u) -- (y);
            \end{tikzpicture}
            \hspace{6ex}
            \begin{tikzpicture}[dgraph]
                \node[cont] (z) at (0,0) {$z$};
                \node[cont] (y) at (1,1) {$y$};
                \node[cont] (x) at (0,2) {$x$};
                \node[cont] (u) at (-1,1) {$u$};
                \draw(x) -- (y);
                \draw(y) -- (z);
                \draw(u) -- (z);
                \draw(x) -- (u);
                \draw(u) -- (y);
            \end{tikzpicture}
        }
    \end{center}
    
    \begin{solution}
        Primeiro, note que ambos os grafos compartilham o mesmo esqueleto e a \'unica raz\~ao pela qual n\~ao est\~ao totalmente conectados \'e a aresta ausente entre $x$ e $z$.
        
        Para o DAG, h\'a tamb\'em apenas uma ordena\'c\~ao que \'e topol\'ogica para o grafo: $x, u, y, z$. A aresta ausente entre $x$ e $y$ corresponde \`a \'unica independ\^encia codificada pelo grafo: $z \independent \pre_z \setminus \pa_z | \pa_z$, ou seja,
        $$z \independent x | u, y.$$ Esta \'e a mesma independ\^encia que obtemos da propriedade de Markov local direcionada.
        
        Para o grafo n\~ao direcionado, 
        $$z \independent x | u, y$$ ocorre porque $u, y$ bloqueiam todos os caminhos entre $z$ e $x$. Todas as vari\'aveis, exceto $z$ e $x$, est\~ao conectadas entre si, de modo que nenhuma outra independ\^encia pode ocorrer.
        
        Portanto, ambos os grafos codificam apenas $z \independent x | u, y$ e s\~ao, assim, I-equivalentes.
        
    \end{solution}
    
    
    \item Os dois grafos a seguir, que s\~ao modelos ocultos de Markov direcionados e n\~ao direcionados, s\~ao I-equivalentes?
    \begin{center}
        \scalebox{0.8}{
            \begin{tikzpicture}[dgraph]
                \node[cont] (y1) at (0,0) {$y_1$};
                \node[cont] (y2) at (2,0) {$y_2$};
                \node[cont] (y3) at (4,0) {$y_3$};
                \node[cont] (y4) at (6,0) {$y_4$};
                \node[cont] (x1) at (0,2) {$x_1$};
                \node[cont] (x2) at (2,2) {$x_2$};
                \node[cont] (x3) at (4,2) {$x_3$};
                \node[cont] (x4) at (6,2) {$x_4$};
                \draw(x1)--(y1);\draw(x2)--(y2);\draw(x3)--(y3);\draw(x4)--(y4);
                \draw(x1)--(x2);\draw(x2)--(x3);\draw(x3)--(x4);
        \end{tikzpicture}}
        \hspace{6ex}
        \scalebox{0.8}{
            \begin{tikzpicture}[ugraph]
                \node[cont] (y1) at (0,0) {$y_1$};
                \node[cont] (y2) at (2,0) {$y_2$};
                \node[cont] (y3) at (4,0) {$y_3$};
                \node[cont] (y4) at (6,0) {$y_4$};
                \node[cont] (x1) at (0,2) {$x_1$};
                \node[cont] (x2) at (2,2) {$x_2$};
                \node[cont] (x3) at (4,2) {$x_3$};
                \node[cont] (x4) at (6,2) {$x_4$};
                \draw(x1)--(y1);\draw(x2)--(y2);\draw(x3)--(y3);\draw(x4)--(y4);
                \draw(x1)--(x2);\draw(x2)--(x3);\draw(x3)--(x4);
        \end{tikzpicture}}
    \end{center}
    
    \begin{solution}
        O esqueleto dos dois grafos \'e o mesmo e n\~ao h\'a imoralidades. Portanto, os dois grafos s\~ao I-equivalentes. 
    \end{solution}
    
    \item Os dois grafos seguintes s\~ao I-equivalentes?
    \begin{center}
        \scalebox{0.8}{
            \begin{tikzpicture}[dgraph]
                \node[cont] (y1) at (0,0) {$y_1$};
                \node[cont] (y2) at (2,0) {$y_2$};
                \node[cont] (y3) at (4,0) {$y_3$};
                \node[cont] (y4) at (6,0) {$y_4$};
                \node[cont] (x1) at (0,2) {$x_1$};
                \node[cont] (x2) at (2,2) {$x_2$};
                \node[cont] (x3) at (4,2) {$x_3$};
                \node[cont] (x4) at (6,2) {$x_4$};
                \draw(x1)--(y1);\draw(x2)--(y2);\draw(x3)--(y3);\draw(x4)--(y4);
                \draw(x1)--(x2);\draw(x3)--(x2);\draw(x3)--(x4);
        \end{tikzpicture}}
        \hspace{6ex}
        \scalebox{0.8}{
            \begin{tikzpicture}[ugraph]
                \node[cont] (y1) at (0,0) {$y_1$};
                \node[cont] (y2) at (2,0) {$y_2$};
                \node[cont] (y3) at (4,0) {$y_3$};
                \node[cont] (y4) at (6,0) {$y_4$};
                \node[cont] (x1) at (0,2) {$x_1$};
                \node[cont] (x2) at (2,2) {$x_2$};
                \node[cont] (x3) at (4,2) {$x_3$};
                \node[cont] (x4) at (6,2) {$x_4$};
                \draw(x1)--(y1);\draw(x2)--(y2);\draw(x3)--(y3);\draw(x4)--(y4);
                \draw(x1)--(x2);\draw(x2)--(x3);\draw(x3)--(x4);
        \end{tikzpicture}}
    \end{center}
    
    \begin{solution}
        Os dois grafos n\~ao s\~ao I-equivalentes porque $x_1-x_2-x_3$ forma uma imoralidade. Assim, o grafo n\~ao direcionado codifica $x_1 \independent x_3 | x_2$, o que n\~ao \'e representado no grafo direcionado. Por outro lado, o grafo direcionado afirma que $x_1 \independent x_3$, o que n\~ao \'e representado no grafo n\~ao direcionado.
    \end{solution}
    
\end{exenumerate}

\ex{Moraliza\'c\~ao: Convertendo DAGs para I-mapas m\'inimos n\~ao direcionados}
\label{ex:DAG-to-undirected}
A seguinte receita constr\'oi I-mapas m\'inimos n\~ao direcionados para $\Ind(p)$:
\begin{itemize}
    \item Determine o Markov blanket para cada vari\'avel $x_i$
    \item Construa um grafo onde os vizinhos de $x_i$ s\~ao dados por seu Markov blanket.
\end{itemize}
Podemos adaptar a receita para construir um I-mapa m\'inimo n\~ao direcionado para as independ\^encias $\Ind(G)$ codificadas por um DAG $G$. O que precisamos fazer \'e usar $G$ para ler os Markov blankets para as vari\'aveis $x_i$, em vez de determinar os Markov blankets a partir da distribui\'c\~ao $p$.

Mostre que esse procedimento leva \`a seguinte receita para converter DAGs em I-mapas m\'inimos n\~ao direcionados:
\begin{enumerate}
    \item Para todas as imoralidades no grafo: adicione arestas entre \emph{todos} os pais do n\'o colisor.
    \item Torne todas as arestas no grafo n\~ao direcionadas.
\end{enumerate}
O primeiro passo \'e s\^o vezes chamado de ``moraliza\'c\~ao'' porque ``casamos'' todos os pais no grafo que ainda n\~ao est\~ao diretamente conectados por uma aresta.

O primeiro passo \'e s\'o vezes chamado de ``moraliza\'c\~ao'' porque n\'os ``casamos'' todos os pais no grafo que ainda n\~ao est\~ao diretamente conectados por uma aresta. O grafo n\~ao direcionado resultante \'e chamado de grafo moral de $G$, s\^o vezes denotado por $\mathcal{M}(G)$.

\begin{solution}
    A cobertura de Markov (Markov blanket) de uma vari\'avel $x$ \'e o conjunto de seus pais, filhos e co-pais, como mostrado no grafo abaixo na subfigura (a). Os pais e filhos est\~ao conectados a $x$ no grafo direcionado, mas os co-pais n\~ao est\~ao diretamente conectados a $x$. Portanto, de acordo com ``Construa um grafo onde os vizinhos de $x_i$ sejam sua cobertura de Markov.'', precisamos introduzir arestas entre $x$ e todos os seus co-pais. Isso resulta no grafo intermedi\'ario na subfigura (b).
    
    Agora, considerando o pai no canto superior esquerdo de $x$, vemos que para aquele n\'o, a cobertura de Markov inclui os outros pais de $x$. Isso significa que precisamos conectar todos os pais de $x$, o que resulta no grafo na subfigura (c). Isso \'e s\^o vezes chamado de ``casar'' os pais de $x$. Continuando dessa forma, vemos que precisamos ``casar'' todos os pais no grafo que ainda n\~ao est\~ao casados.
    
    Finalmente, precisamos tornar todas as arestas no grafo n\~ao direcionadas, o que resulta na subfigura (d).
    
    Uma abordagem mais simples \'e notar que o DAG especifica a fatoriza\'c\~ao $p(\x) = \prod_i p(x_i | \pa_i)$. Podemos considerar cada condicional $p(x_i | \pa_i)$ como um fator $\phi_i(x_i, \pa_i)$ de modo que obtemos a distribui\'c\~ao de Gibbs $p(\x) = \prod_i \phi_i(x_i | \pa_i)$. Visualizar a distribui\'c\~ao conectando todas as vari\'aveis no mesmo fator $\phi_i(x_i | \pa_i)$ leva ao ``casamento'' de todos os pais de $x_i$. Isso corresponde ao primeiro passo na receita porque $x_i$ est\'a em uma configura\'c\~ao de colisor em rela\'c\~ao aos n\'os pais. Nem todos os pais formam uma imoralidade, mas isso aqui n\~ao importa porque aqueles que n\~ao formam uma imoralidade j\'a est\~ao conectados por uma aresta de cobertura em primeiro lugar.
    
    \begin{figure}[htb]
        \centering
        \begin{tabular}[b]{c}
            \scalebox{0.7}{ % x,y
                \begin{tikzpicture}[dgraph]
                    % x,y
                    
                    \node[cont] (lowest1) at ( 2,0) {};
                    \node[cont] (lowest2) at ( 3.5,0) {};
                    \node[cont] (lowest3) at ( 4.5,0) {};
                    
                    \node[cont] (belowleft) at ( 2,1) {};
                    \node[cont] (belowright) at ( 4,1) {};
                    
                    %\node[cont] at ( 0,2) {};
                    \node[cont] (copleft) at ( 1,2) {};
                    %\node[cont] at ( 2,2) {};
                    \node[cont] (x) at ( 3,2) {$x$};
                    % \node[cont] at ( 4,2) {};
                    \node[cont] (copright) at ( 5,2) {};
                    
                    \node (abovemostleft) at ( 1,3) {};
                    \node[cont] (aboveleft) at ( 2,3) {};
                    \node[cont] (above) at ( 3,3) {};
                    \node[cont] (aboveright) at ( 4,3) {};
                    \node (abovemostright) at ( 5,3) {};
                    
                    \node (top1) at ( 1,4) {};
                    \node[cont] (top2) at ( 2,4) {};
                    \node (top3) at ( 3,4) {};
                    \node[cont] (top4) at ( 4,4) {};
                    \node (top5) at ( 5,4) {};
                    
                    \draw (top1) -- (aboveleft);
                    \draw (top2) -- (aboveleft);
                    \draw (top3) -- (above);
                    \draw (top4) -- (aboveright);
                    \draw (top5) -- (aboveright);
                    
                    \draw (aboveright) -- (x);
                    \draw (aboveleft) -- (x);
                    \draw (above) -- (x);
                    
                    \draw (x) -- (belowright);
                    \draw (x) -- (belowleft);
                    
                    \draw (copleft) -- (belowleft);
                    \draw (copright) -- (belowright);
                    
                    \draw (abovemostleft) -- (copleft);
                    \draw (abovemostright) -- (copright);
                    
                    \draw (belowright) -- (lowest2);
                    \draw (belowright) -- (lowest3);
                    \draw (belowleft) -- (lowest1);
                    
                    \node[contobs] at ( 2,3) {};
                    \node[contobs] at ( 3,3) {};
                    \node[contobs] at ( 4,3) {};
                    
                    \node[contobs] at ( 2,1) {};
                    \node[contobs] at ( 4,1) {};
                    
                    \node[contobs] at ( 1,2) {};
                    \node[contobs] at ( 5,2) {};
            \end{tikzpicture}}\\
            {\small (a) DAG}
        \end{tabular}
        \begin{tabular}[b]{c}
            \scalebox{0.7}{ % x,y
                \begin{tikzpicture}[dgraph]
                    % x,y
                    
                    \node[cont] (lowest1) at ( 2,0) {};
                    \node[cont] (lowest2) at ( 3.5,0) {};
                    \node[cont] (lowest3) at ( 4.5,0) {};
                    
                    \node[cont] (belowleft) at ( 2,1) {};
                    \node[cont] (belowright) at ( 4,1) {};
                    
                    %\node[cont] at ( 0,2) {};
                    \node[cont] (copleft) at ( 1,2) {};
                    %\node[cont] at ( 2,2) {};
                    \node[cont] (x) at ( 3,2) {$x$};
                    % \node[cont] at ( 4,2) {};
                    \node[cont] (copright) at ( 5,2) {};
                    
                    \node (abovemostleft) at ( 1,3) {};
                    \node[cont] (aboveleft) at ( 2,3) {};
                    \node[cont] (above) at ( 3,3) {};
                    \node[cont] (aboveright) at ( 4,3) {};
                    \node (abovemostright) at ( 5,3) {};
                    
                    \node (top1) at ( 1,4) {};
                    \node[cont] (top2) at ( 2,4) {};
                    \node (top3) at ( 3,4) {};
                    \node[cont] (top4) at ( 4,4) {};
                    \node (top5) at ( 5,4) {};
                    
                    \draw (top1) -- (aboveleft);
                    \draw (top2) -- (aboveleft);
                    \draw (top3) -- (above);
                    \draw (top4) -- (aboveright);
                    \draw (top5) -- (aboveright);
                    
                    \draw (aboveright) -- (x);
                    \draw (aboveleft) -- (x);
                    \draw (above) -- (x);
                    
                    \draw (x) -- (belowright);
                    \draw (x) -- (belowleft);
                    
                    \draw (copleft) -- (belowleft);
                    \draw (copright) -- (belowright);
                    
                    \draw (abovemostleft) -- (copleft);
                    \draw (abovemostright) -- (copright);
                    
                    \draw (belowright) -- (lowest2);
                    \draw (belowright) -- (lowest3);
                    \draw (belowleft) -- (lowest1);
                    
                    \node[contobs] at ( 2,3) {};
                    \node[contobs] at ( 3,3) {};
                    \node[contobs] at ( 4,3) {};
                    
                    \node[contobs] at ( 2,1) {};
                    \node[contobs] at ( 4,1) {};
                    
                    \node[contobs] at ( 1,2) {};
                    \node[contobs] at ( 5,2) {};
                    
                    \draw[-] (copleft) -- (x);
                    \draw[-] (copright) -- (x);
            \end{tikzpicture}}\\
            {\small  (b) Passo intermedi\'ario 1}
        \end{tabular}
        \begin{tabular}[b]{c}
            \scalebox{0.7}{ % x,y
                \begin{tikzpicture}[dgraph]
                    % x,y
                    
                    \node[cont] (lowest1) at ( 2,0) {};
                    \node[cont] (lowest2) at ( 3.5,0) {};
                    \node[cont] (lowest3) at ( 4.5,0) {};
                    
                    \node[cont] (belowleft) at ( 2,1) {};
                    \node[cont] (belowright) at ( 4,1) {};
                    
                    %\node[cont] at ( 0,2) {};
                    \node[cont] (copleft) at ( 1,2) {};
                    %\node[cont] at ( 2,2) {};
                    \node[cont] (x) at ( 3,2) {$x$};
                    % \node[cont] at ( 4,2) {};
                    \node[cont] (copright) at ( 5,2) {};
                    
                    \node (abovemostleft) at ( 1,3) {};
                    \node[cont] (aboveleft) at ( 2,3) {};
                    \node[cont] (above) at ( 3,3) {};
                    \node[cont] (aboveright) at ( 4,3) {};
                    \node (abovemostright) at ( 5,3) {};
                    
                    \node (top1) at ( 1,4) {};
                    \node[cont] (top2) at ( 2,4) {};
                    \node (top3) at ( 3,4) {};
                    \node[cont] (top4) at ( 4,4) {};
                    \node (top5) at ( 5,4) {};
                    
                    \draw (top1) -- (aboveleft);
                    \draw (top2) -- (aboveleft);
                    \draw (top3) -- (above);
                    \draw (top4) -- (aboveright);
                    \draw (top5) -- (aboveright);
                    
                    \draw (aboveright) -- (x);
                    \draw (aboveleft) -- (x);
                    \draw (above) -- (x);
                    
                    \draw (x) -- (belowright);
                    \draw (x) -- (belowleft);
                    
                    \draw (copleft) -- (belowleft);
                    \draw (copright) -- (belowright);
                    
                    \draw (abovemostleft) -- (copleft);
                    \draw (abovemostright) -- (copright);
                    
                    \draw (belowright) -- (lowest2);
                    \draw (belowright) -- (lowest3);
                    \draw (belowleft) -- (lowest1);
                    
                    \node[contobs] at ( 2,3) {};
                    \node[contobs] at ( 3,3) {};
                    \node[contobs] at ( 4,3) {};
                    
                    \node[contobs] at ( 2,1) {};
                    \node[contobs] at ( 4,1) {};
                    
                    \node[contobs] at ( 1,2) {};
                    \node[contobs] at ( 5,2) {};
                    
                    \draw[-] (copleft) -- (x);
                    \draw[-] (copright) -- (x);
                    \draw[-] (aboveleft) -- (above);
                    \draw[-] (aboveleft) to [bend right=-45] (aboveright);
                    \draw[-] (above) -- (aboveright);
            \end{tikzpicture}}\\
            {\small (c) Passo intermedi\'ario 2}
        \end{tabular}
        \begin{tabular}[b]{c}
            \scalebox{0.7}{ % x,y
                \begin{tikzpicture}[ugraph]
                    % x,y
                    
                    \node[cont] (lowest1) at ( 2,0) {};
                    \node[cont] (lowest2) at ( 3.5,0) {};
                    \node[cont] (lowest3) at ( 4.5,0) {};
                    
                    \node[cont] (belowleft) at ( 2,1) {};
                    \node[cont] (belowright) at ( 4,1) {};
                    
                    %\node[cont] at ( 0,2) {};
                    \node[cont] (copleft) at ( 1,2) {};
                    %\node[cont] at ( 2,2) {};
                    \node[cont] (x) at ( 3,2) {$x$};
                    % \node[cont] at ( 4,2) {};
                    \node[cont] (copright) at ( 5,2) {};
                    
                    \node (abovemostleft) at ( 1,3) {};
                    \node[cont] (aboveleft) at ( 2,3) {};
                    \node[cont] (above) at ( 3,3) {};
                    \node[cont] (aboveright) at ( 4,3) {};
                    \node (abovemostright) at ( 5,3) {};
                    
                    \node (top1) at ( 1,4) {};
                    \node[cont] (top2) at ( 2,4) {};
                    \node (top3) at ( 3,4) {};
                    \node[cont] (top4) at ( 4,4) {};
                    \node (top5) at ( 5,4) {};
                    
                    \draw (top1) -- (aboveleft);
                    \draw (top2) -- (aboveleft);
                    \draw (top3) -- (above);
                    \draw (top4) -- (aboveright);
                    \draw (top5) -- (aboveright);
                    
                    \draw (aboveright) -- (x);
                    \draw (aboveleft) -- (x);
                    \draw (above) -- (x);
                    
                    \draw (x) -- (belowright);
                    \draw (x) -- (belowleft);
                    
                    \draw (copleft) -- (belowleft);
                    \draw (copright) -- (belowright);
                    
                    \draw (abovemostleft) -- (copleft);
                    \draw (abovemostright) -- (copright);
                    
                    \draw (belowright) -- (lowest2);
                    \draw (belowright) -- (lowest3);
                    \draw (belowleft) -- (lowest1);
                    
                    \node[contobs] at ( 2,3) {};
                    \node[contobs] at ( 3,3) {};
                    \node[contobs] at ( 4,3) {};
                    
                    \node[contobs] at ( 2,1) {};
                    \node[contobs] at ( 4,1) {};
                    
                    \node[contobs] at ( 1,2) {};
                    \node[contobs] at ( 5,2) {};
                    
                    \draw[-] (copleft) -- (x);
                    \draw[-] (copright) -- (x);
                    \draw[-] (aboveleft) -- (above);
                    \draw[-] (aboveleft) to [bend right=-45] (aboveright);
                    \draw[-] (above) -- (aboveright);
                    \draw[-] (top1) to [bend right=-45] (top2);
                    \draw[-] (top4) to [bend right=-45] (top5); 
            \end{tikzpicture}}\\
            {\small (d) Grafo n\~ao direcionado}
        \end{tabular}
        \caption{Resposta ao \exref{ex:DAG-to-undirected}: Ilustrando o processo de moraliza\'c\~ao}
    \end{figure}
\end{solution}

\ex{Exerc\'icio de moraliza\'c\~ao}
\label{ex:moralisation-alt}

Para o DAG $G$ abaixo, encontre o I-map n\~ao direcionado m\'inimo para $\Ind(G)$.

\begin{center}
    \scalebox{0.8}{
        \begin{tikzpicture}[dgraph]
            
            \node[cont] (x2) at (2,0) {$x_2$};
            \node[cont, left = of x2] (x1) {$x_1$};
            \node[cont, right = of x2] (x3) {$x_3$};
            \node[cont, below = of x2] (x4) {$x_4$};
            \node[cont, right = of x4] (x5) {$x_5$};
            \node[cont, below = of x4] (x6) {$x_6$};
            \node[cont, below = of x5] (x7) {$x_7$};
            
            \draw (x1) -- (x4);
            \draw (x2) -- (x4);
            \draw (x3) -- (x4);
            \draw (x4) -- (x6);
            \draw (x4) -- (x7);
            \draw (x5) -- (x7);
            
        \end{tikzpicture}
    }
\end{center}

\begin{solution}
    
    Para derivar um I-map n\~ao direcionado m\'inimo a partir de um direcionado, temos que construir o grafo moralizado onde os pais ``n\~ao casados'' s\~ao conectados por uma aresta de cobertura. Isso ocorre porque cada condicional $p(x_i | \pa_i)$ corresponde a um fator $\phi_i(x_i,\pa_i)$ e precisamos conectar todas as vari\'aveis que s\~ao argumentos do mesmo fator com arestas.
    
    Estatisticamente, a raz\~ao para casar os pais \'e a seguinte: Uma independ\^encia $x \independent y | \{\text{filho, outros n\'os}\}$ n\~ao se sustenta no grafo direcionado em caso de conex\~oes de colisor, mas se sustentaria no grafo n\~ao direcionado se n\~ao cas\'assemos os pais. Portanto, links entre os pais devem ser adicionados.
    
    \'E importante adicionar arestas entre \emph{todos} os pais de um n\'o. Aqui, $p(x_4 | x_1, x_2, x_3)$ corresponde a um fator $\phi(x_4, x_1, x_2, x_3)$ de modo que todas as quatro vari\'aveis precisam estar conectadas. Apenas adicionar as arestas $x_1 - x_2$ e $x_2 - x_3$ n\~ao seria suficiente.
    
    O grafo moral, que \'e o I-map n\~ao direcionado m\'inimo solicitado, \'e mostrado abaixo.
    
    \begin{center}
        \scalebox{0.8}{
            \begin{tikzpicture}[ugraph]
                
                \node[cont] (x2) at (2,0) {$x_2$};
                \node[cont, left = of x2] (x1) {$x_1$};
                \node[cont, right = of x2] (x3) {$x_3$};
                \node[cont, below = of x2] (x4) {$x_4$};
                \node[cont, right = of x4] (x5) {$x_5$};
                \node[cont, below = of x4] (x6) {$x_6$};
                \node[cont, below = of x5] (x7) {$x_7$};
                
                \draw (x1) -- (x4);
                \draw (x2) -- (x4);
                \draw (x3) -- (x4);
                \draw (x4) -- (x6);
                \draw (x4) -- (x7);
                \draw (x5) -- (x7);
                
                %arestas de cobertura
                \draw (x1) -- (x2);
                \draw (x2) -- (x3);
                \draw (x1) to [bend left=45] (x3);
                \draw (x4) -- (x5);
                
            \end{tikzpicture}
        }
    \end{center}
    
\end{solution}

% -------------------------------------------------------------------
\ex{Exerc\'icio de moraliza\'c\~ao}

Considere o DAG $G$:
\begin{center}
    \scalebox{0.9}{ % x,y
        \begin{tikzpicture}[dgraph]
            \node[cont] (y) at (0,0) {$y$};
            \node[cont] (z1) at (-1.75,1.5) {$z_1$};
            \node[cont] (z2) at (1.75,1.5) {$z_2$};
            \node[cont] (x1) at (-3,3) {$x_1$};
            \node[cont] (x2) at (-1.75,3) {$x_2$};
            \node[cont] (x3) at (-0.5,3) {$x_3$};
            \node[cont] (x4) at (0.5,3) {$x_4$};
            \node[cont] (x5) at (1.75,3) {$x_5$};
            \node[cont] (x6) at (3,3) {$x_6$};
            
            \draw (z1) -- (y);
            \draw (z2) -- (y);
            \draw (x1) -- (z1);
            \draw (x2) -- (z1);
            \draw (x3) -- (z1);
            \draw (x4) -- (z2);
            \draw (x5) -- (z2);
            \draw (x6) -- (z2);
    \end{tikzpicture}}
\end{center}
Um amigo afirma que o grafo n\~ao direcionado abaixo \'e o grafo moral $\mathcal{M}(G)$ de $G$. Seu amigo est\'a correto? Se n\~ao, indique quais arestas precisariam ser removidas ou adicionadas e explique, em termos de independ\^encias representadas, por que as altera\'c\~oes s\~ao necess\'arias para que o grafo se torne o grafo moral de $G$. 

\begin{center}
    \scalebox{0.9}{ % x,y
        \begin{tikzpicture}[ugraph]
            \node[cont] (y) at (0,0) {$y$};
            \node[cont] (z1) at (-1.75,1.5) {$z_1$};
            \node[cont] (z2) at (1.75,1.5) {$z_2$};
            \node[cont] (x1) at (-3,3) {$x_1$};
            \node[cont] (x2) at (-1.75,3) {$x_2$};
            \node[cont] (x3) at (-0.5,3) {$x_3$};
            \node[cont] (x4) at (0.5,3) {$x_4$};
            \node[cont] (x5) at (1.75,3) {$x_5$};
            \node[cont] (x6) at (3,3) {$x_6$};
            
            \draw (z1) -- (y);
            \draw (z2) -- (y);
            \draw (x1) -- (z1);
            \draw (x2) -- (z1);
            \draw (x3) -- (z1);
            \draw (x4) -- (z2);
            \draw (x5) -- (z2);
            \draw (x6) -- (z2);
            
            \draw (x1) -- (x2);
            \draw (x2) -- (x3);
            \draw (x4) -- (x5);
            \draw (x5) -- (x6);
            
            \draw (x1) to [out=40,in=140] (x6);
            
            \draw (z1) -- (z2);
    \end{tikzpicture}}
\end{center}

\begin{solution}
    O grafo moral $\mathcal{M}(G)$ \'e um I-map n\~ao direcionado m\'inimo das independ\^encias representadas por $G$. Seguindo o procedimento de conectar pais ``n\~ao casados'' de colisores, obtemos o seguinte grafo moral de $G$:
    
    \begin{center}
        \begin{tikzpicture}[ugraph]
            \node[cont] (y) at (0,0) {$y$};
            \node[cont] (z1) at (-1.75,1.5) {$z_1$};
            \node[cont] (z2) at (1.75,1.5) {$z_2$};
            \node[cont] (x1) at (-3,3) {$x_1$};
            \node[cont] (x2) at (-1.75,3) {$x_2$};
            \node[cont] (x3) at (-0.5,3) {$x_3$};
            \node[cont] (x4) at (0.5,3) {$x_4$};
            \node[cont] (x5) at (1.75,3) {$x_5$};
            \node[cont] (x6) at (3,3) {$x_6$};
            
            \draw (z1) -- (y);
            \draw (z2) -- (y);
            \draw (x1) -- (z1);
            \draw (x2) -- (z1);
            \draw (x3) -- (z1);
            \draw (x4) -- (z2);
            \draw (x5) -- (z2);
            \draw (x6) -- (z2);
            
            \draw (x1) -- (x2);
            \draw (x2) -- (x3);
            \draw (x1) to [out=45,in=135] (x3);
            
            \draw (x4) -- (x5);
            \draw (x5) -- (x6);
            \draw (x4) to [out=45,in=135] (x6);
            
            \draw (z1) -- (z2);
        \end{tikzpicture}
        %}
\end{center}
Podemos assim ver que o grafo n\~ao direcionado do amigo n\~ao \'e o grafo moral de $G$.

A aresta entre $x_1$ e $x_6$ pode ser removida. Isso ocorre porque para $G$, temos por exemplo as independ\^encias $x_1 \independent x_6 | z_1$, $x_1 \independent x_6 | z_2$, $x_1 \independent x_6 | z_1, z_2$, as quais n\~ao s\~ao representadas pelo grafo n\~ao direcionado desenhado.

Precisamos adicionar arestas entre $x_1$ e $x_3$, e entre $x_4$ e $x_6$. Caso contr\'ario, o grafo n\~ao direcionado faz a asser\'c\~ao de independ\^encia errada de que $x_1 \independent x_3 | x_2, z_1$ (e o equivalente para $x_4$ e $x_6$).

\end{solution}


\ex{Triangula\'c\~ao: Convertendo grafos n\~ao direcionados em I-maps direcionados m\'inimos}
\label{ex:undirected-to-DAG-alt}

No \exref{ex:DAG-to-undirected} adaptamos uma receita para construir I-maps n\~ao direcionados m\'inimos para $\Ind(p)$ para o caso de $\Ind(G)$, onde $G$ \'e um DAG. A diferen\'ca fundamental foi que usamos o grafo $G$ para determinar independ\^encias em vez da distribui\'c\~ao $p$.

Podemos adaptar de forma semelhante a receita para construir um I-map direcionado m\'inimo para $\Ind(p)$ para construir um I-map direcionado m\'inimo para $\Ind(H)$, onde $H$ \'e um grafo n\~ao direcionado:
\begin{enumerate}
\item Escolha uma ordena\'c\~ao das vari\'aveis aleat\'orias.
\item Para todas as vari\'aveis $x_i$, use $H$ para determinar um subconjunto \emph{m\'inimo} $\pap_i$ dos predecessores $\pre_i$ tal que
$$ x_i \independent \left(\pre_i \setminus \pap_i\right) \mid \pap_i $$
se sustenta.
\item Construa um DAG com os $\pap_i$ como pais $\pa_i$ de $x_i$.
\end{enumerate}
Observa\'c\~oes: (1) I-maps direcionados m\'inimos obtidos com ordena\'c\~oes diferentes n\~ao s\~ao geralmente I-equivalentes. (2) Os I-maps direcionados m\'inimos obtidos com o m\'etodo acima s\~ao sempre grafos cordais (chordal graphs). Grafos cordais s\~ao grafos onde a trilha mais longa sem atalhos \'e um tri\^angulo (\url{https://en.wikipedia.org/wiki/Chordal_graph}). Eles s\~ao, portanto, tamb\'em chamados de grafos triangulados. Obtemos grafos cordais porque se tiv\'essemos trilhas sem atalhos que envolvessem mais de 3 n\'os, ter\'iamos necessariamente uma imoralidade no grafo. Mas imoralidades codificam independ\^encias que um grafo n\~ao direcionado n\~ao pode representar, o que faria com que o DAG n\~ao fosse mais um I-map para $\Ind(H)$.

\begin{exenumerate}
\item Seja $H$ o grafo n\~ao direcionado abaixo. Determine o I-map direcionado m\'inimo para $\Ind(H)$ com a ordena\'c\~ao de vari\'aveis $x_1, x_2, x_3, x_4, x_5$.
\label{q:undirected-to-directed-I-map-alt}
\begin{center}
    \scalebox{0.9}{ % x,y
        \begin{tikzpicture}[ugraph]
            \node[cont] (x1) at (0,0) {$x_1$};
            \node[cont] (x2) at (2,0) {$x_2$};
            \node[cont] (x3) at (0,-1.5) {$x_3$};
            \node[cont] (x4) at (2,-1.5) {$x_4$};
            \node[cont] (x5) at (1,-3) {$x_5$};
            
            \draw(x1) -- (x2);
            \draw(x1) -- (x3);
            \draw(x3) -- (x5);
            \draw(x2) -- (x4);
            \draw(x4) -- (x5);
    \end{tikzpicture}}
\end{center}

\begin{solution}
    Usamos a ordena\'c\~ao $x_1, x_2, x_3, x_4, x_5$ e seguimos o procedimento de convers\~ao:
    
    \begin{itemize}
        \item $x_2$ n\~ao \'e independente de $x_1$, de modo que definimos $\pa_2 =\{x_1\}$. Veja o primeiro grafo na Figura \ref{fig:undirected-to-directed-I-map-alt}.
        \item Como $x_3$ est\'a conectado tanto a $x_1$ quanto a $x_2$, n\~ao temos $x_3 \independent x_2, x_1$. N\~ao podemos tornar $x_3$ independente de $x_2$ condicionando em $x_1$ porque existem dois caminhos de $x_3$ para $x_2$ e $x_1$ apenas bloqueia o superior. Al\'em disso, $x_1$ \'e um vizinho de $x_3$, de modo que condicionar em $x_2$ n\~ao os torna independentes. Portanto, devemos definir $\pa_3 = \{x_1, x_2\}$. Veja o segundo grafo na Figura \ref{fig:undirected-to-directed-I-map-alt}.
        \item Para $x_4$, vemos a partir do grafo n\~ao direcionado que $x_4 \independent x_1 \mid x_3, x_2$. O grafo mostra ainda que remover $x_3$ ou $x_2$ do conjunto de condicionamento n\~ao \'e poss\'ivel, e condicionar em $x_1$ n\~ao tornar\'a $x_4$ independente de $x_2$ ou $x_3$. Temos assim $\pa_4 = \{x_2, x_3\}$. Veja o quarto grafo na Figura \ref{fig:undirected-to-directed-I-map-alt}.
        \item O mesmo racioc\'inio mostra que $\pa_5 = \{x_3,x_4\}$. Veja o \'ultimo grafo na Figura \ref{fig:undirected-to-directed-I-map-alt}.
    \end{itemize}
    Isso resulta no grafo direcionado triangulado na Figura \ref{fig:undirected-to-directed-I-map-alt} \`a direita.
    
    \begin{figure}[h!]
        \centering
        \scalebox{0.9}{ % x,y
            \begin{tikzpicture}[dgraph]
                \node[cont] (x1) at (0,0) {$x_1$};
                \node[cont] (x2) at (2,0) {$x_2$};
                \node[cont] (x3) at (0,-1.5) {$x_3$};
                \node[cont] (x4) at (2,-1.5) {$x_4$};
                \node[cont] (x5) at (1,-3) {$x_5$};
                
                \draw(x1) -- (x2);
                
                \draw[-,gray,dotted] (2.5,-3) -- (2.5,0);
            \end{tikzpicture}
            \hspace{2ex}
            \begin{tikzpicture}[dgraph]
                \node[cont] (x1) at (0,0) {$x_1$};
                \node[cont] (x2) at (2,0) {$x_2$};
                \node[cont] (x3) at (0,-1.5) {$x_3$};
                \node[cont] (x4) at (2,-1.5) {$x_4$};
                \node[cont] (x5) at (1,-3) {$x_5$};
                
                \draw(x1) -- (x2);
                \draw(x1) -- (x3);
                \draw(x2) -- (x3);
                
                \draw[-,gray,dotted] (2.5,-3) -- (2.5,0);
            \end{tikzpicture}
            \hspace{2ex}
            \begin{tikzpicture}[dgraph]
                \node[cont] (x1) at (0,0) {$x_1$};
                \node[cont] (x2) at (2,0) {$x_2$};
                \node[cont] (x3) at (0,-1.5) {$x_3$};
                \node[cont] (x4) at (2,-1.5) {$x_4$};
                \node[cont] (x5) at (1,-3) {$x_5$};
                
                \draw(x1) -- (x2);
                \draw(x1) -- (x3);
                \draw(x2) -- (x3);
                \draw(x2) -- (x4);
                \draw(x3) -- (x4);
                
                \draw[-,gray,dotted] (2.5,-3) -- (2.5,0);
            \end{tikzpicture}
            \hspace{2ex}
            \begin{tikzpicture}[dgraph]
                \node[cont] (x1) at (0,0) {$x_1$};
                \node[cont] (x2) at (2,0) {$x_2$};
                \node[cont] (x3) at (0,-1.5) {$x_3$};
                \node[cont] (x4) at (2,-1.5) {$x_4$};
                \node[cont] (x5) at (1,-3) {$x_5$};
                
                \draw(x1) -- (x2);
                \draw(x1) -- (x3);
                \draw(x2) -- (x3);
                \draw(x2) -- (x4);
                \draw(x3) -- (x4);
                \draw(x3) -- (x5);
                \draw(x4) -- (x5);
            \end{tikzpicture}
        }
        \caption{\label{fig:undirected-to-directed-I-map-alt}. Resposta ao \exref{ex:undirected-to-DAG-alt}, Quest\~ao \ref{q:undirected-to-directed-I-map-alt}. }
        
    \end{figure}
    
    Para ver por que a triangula\'c\~ao \'e necess\'aria, considere o caso em que n\~ao tiv\'essemos a aresta entre $x_2$ e $x_3$, como na Figura \ref{fig:triangulation-alt}. O grafo direcionado ent\~ao implicaria que $x_3 \independent x_2 \mid x_1$ (verifique!). Mas esta asser\'c\~ao de independ\^encia n\~ao se sustenta no grafo n\~ao direcionado, de modo que o grafo na Figura \ref{fig:triangulation-alt} n\~ao \'e um I-map.
    
    \begin{figure}[h!]
        \centering
        \begin{tikzpicture}[dgraph]
            \node[cont] (x1) at (0,0) {$x_1$};
            \node[cont] (x2) at (2,0) {$x_2$};
            \node[cont] (x3) at (0,-1.5) {$x_3$};
            \node[cont] (x4) at (2,-1.5) {$x_4$};
            \node[cont] (x5) at (1,-3) {$x_5$};
            
            \draw(x1) -- (x2);
            \draw(x1) -- (x3);
            % \draw(x2) -- (x3);
            \draw(x2) -- (x4);
            \draw(x3) -- (x4);
            \draw(x3) -- (x5);
            \draw(x4) -- (x5);
        \end{tikzpicture}
        \caption{\label{fig:triangulation-alt} N\~ao \'e um I-map direcionado para o modelo gr\'afico n\~ao direcionado definido pelo grafo no \exref{ex:undirected-to-DAG-alt}, Quest\~ao \ref{q:undirected-to-directed-I-map-alt}.}
    \end{figure}
\end{solution}

\item Para o grafo n\~ao direcionado da quest\~ao \ref{q:undirected-to-directed-I-map-alt} acima, qual ordena\'c\~ao de vari\'aveis resulta no I-map direcionado m\'inimo abaixo?
\label{q:undirected-to-directed-I-map-chordal-graph-variable-ordering-alt}
\begin{center}
    \begin{tikzpicture}[dgraph]
        \node[cont] (x1) at (0,0) {$x_1$};
        \node[cont] (x2) at (2,0) {$x_2$};
        \node[cont] (x3) at (0,-1.5) {$x_3$};
        \node[cont] (x4) at (2,-1.5) {$x_4$};
        \node[cont] (x5) at (1,-3) {$x_5$};
        
        \draw(x1) -- (x2);
        \draw(x1) -- (x3);
        \draw(x1) -- (x4);
        \draw(x2) -- (x4);
        \draw(x4) -- (x3);
        \draw(x3) -- (x5);
        \draw(x4) -- (x5);
    \end{tikzpicture}
\end{center}

\begin{solution}
    $x_1$ \'e a raiz do DAG, portanto ele vem primeiro. Em seguida na ordena\'c\~ao est\~ao os filhos de $x_1$: $x_2, x_3, x_4$. Como $x_3$ \'e filho de $x_4$, e $x_4$ um filho de $x_2$, devemos ter $x_1, x_2, x_4, x_3$. Al\'em disso, $x_3$ deve vir antes de $x_5$ na ordena\'c\~ao, j\'a que $x_5$ \'e um filho de $x_3$, logo a ordena\'c\~ao usada deve ter sido: $x_1, x_2, x_4, x_3, x_5$.
    
\end{solution}

\end{exenumerate}


\ex{I-maps, I-maps m\'inimos e I-equival\^encia}
Considere a seguinte fun\'c\~ao de densidade de probabilidade para as vari\'aveis aleat\'orias $x_1, \ldots, x_6$.
$$ p_a(x_1, \ldots, x_6) = p(x_1) p(x_2) p(x_3 | x_1,x_2) p(x_4 | x_2)
p(x_5| x_1) p(x_6 | x_3, x_4, x_5)$$ Para cada um dos dois grafos abaixo, explique se ele \'e um I-map m\'inimo, n\~ao um I-map m\'inimo mas ainda assim um I-map, ou se n\~ao \'e um I-map para as independ\^encias que valem para $p_a$.

\begin{center}
\scalebox{0.95}{ % x,y
    \begin{tikzpicture}[dgraph]
        \node[cont] (x1) at (0,0) {$x_1$};
        \node[cont] (x2) at (2,0) {$x_2$};
        \node[cont] (x3) at (1,-1.5) {$x_3$};
        \node[cont] (x4) at (3,-1.5) {$x_4$};
        \node[cont] (x5) at (-1,-1.5) {$x_5$};
        \node[cont] (x6) at (1,-3) {$x_6$};
        
        \draw (x1) -- (x3);
        \draw (x2) -- (x3);
        \draw (x2) -- (x4);
        \draw (x1) -- (x5);
        \draw (x3) -- (x6);
        \draw (x5) -- (x6);
        \draw (x4) -- (x6);
        
        \draw (x3) -- (x4);
        \draw (x3) -- (x5);
        
        \node[] (text) at (1,-4) {grafo 1};
    \end{tikzpicture}\hspace{12ex}
    \begin{tikzpicture}[ugraph]
        \node[cont] (x1) at (0,0) {$x_1$};
        \node[cont] (x2) at (2,0) {$x_2$};
        \node[cont] (x3) at (1,-1.5) {$x_3$};
        \node[cont] (x4) at (3,-1.5) {$x_4$};
        \node[cont] (x5) at (-1,-1.5) {$x_5$};
        \node[cont] (x6) at (1,-3) {$x_6$};
        
        \draw (x1) -- (x3);
        \draw (x2) -- (x3);
        \draw (x2) -- (x4);
        \draw (x1) -- (x5);
        \draw (x3) -- (x6);
        \draw (x5) -- (x6);
        \draw (x4) -- (x6);
        
        \draw (x3) -- (x4);
        \draw (x3) -- (x5);
        \draw (x1) -- (x2);
        \node[] (text) at (1,-4) {grafo 2};
\end{tikzpicture}}
\end{center}

\begin{solution}
A fdp pode ser visualizada como o seguinte grafo direcionado, o qual \'e um I-map m\'inimo para ela.\\
\begin{center}
    \scalebox{1}{ % x,y
        \begin{tikzpicture}[dgraph]
            \node[cont] (x1) at (0,0) {$x_1$};
            \node[cont] (x2) at (2,0) {$x_2$};
            \node[cont] (x3) at (1,-1.5) {$x_3$};
            \node[cont] (x4) at (3,-1.5) {$x_4$};
            \node[cont] (x5) at (-1,-1.5) {$x_5$};
            \node[cont] (x6) at (1,-3) {$x_6$};
            
            \draw (x1) -- (x3);
            \draw (x2) -- (x3);
            \draw (x2) -- (x4);
            \draw (x1) -- (x5);
            \draw (x3) -- (x6);
            \draw (x5) -- (x6);
            \draw (x4) -- (x6);
            
    \end{tikzpicture}}
\end{center}
O Grafo 1 define distribui\'c\~oes que se fatorizam como
\begin{equation}
    p_b(\x) = p(x_1) p(x_2) p(x_3|x_1, x_2) p(x_4|x_2, x_3) p(x_5|x_1, x_3) p(x_6|x_3,x_4,x_5).
\end{equation}
Comparando com $p_a(x_1, \ldots, x_6)$, vemos que apenas os condicionais $p(x_4|x_2, x_3)$ e $p(x_5|x_1, x_3)$ s\~ao diferentes. Especificamente, seu conjunto de condicionamento inclui $x_3$, o que significa que o Grafo 1 codifica menos independ\^encias do que o que $p_a(x_1, \ldots, x_6)$ satisfaz. Em particular, $x_4 \independent x_3 | x_2$ e $x_5 \independent x_3 | x_1$ n\~ao est\~ao representados no grafo. Isso significa que poder\'iamos remover $x_3$ dos conjuntos de condicionamento, ou equivalentemente remover as arestas $x_3 \rightarrow x_4$ e $x_3 \rightarrow x_5$ do grafo sem introduzir asser\'c\~oes de independ\^encia que n\~ao se sustentem para $p_a$. Isso significa que o grafo 1 \'e um I-map, mas n\~ao um I-map m\'inimo.

O Grafo 2 n\~ao \'e um I-map. Para ser um I-map n\~ao direcionado m\'inimo, ter\'iamos que conectar as vari\'aveis $x_5$ e $x_4$, que s\~ao pais de $x_6$. O Grafo 2 afirma erroneamente que $x_5 \independent x_4 \mid x_1,x_3,x_6$.

\end{solution}


\ex{Limites de modelos gr\'aficos direcionados e n\~ao direcionados}

Consideramos aqui o modelo probabil\'istico $p(y_1,y_2,x_1,x_2) = p(y_1, y_2 | x_1, x_2)p(x_1) p(x_2)$ onde $p(y_1, y_2 | x_1, x_2)$ se fatoriza como
\begin{equation}
p(y_1, y_2 | x_1, x_2) = p(y_1 | x_1) p(y_2 | x_2) \phi(y_1,y_2) n(x_1,x_2)
\end{equation}
com $n(x_1,x_2)$ igual a
\begin{equation}
n(x_1,x_2) = \left(\int p(y_1 | x_1) p(y_2 | x_2) \phi(y_1,y_2) \ud y_1 \ud y_2\right)^{-1}.
\label{eq:ndef-alt}
\end{equation}
No modelo, $x_1$ e $x_2$ s\~ao duas entradas independentes que controlam, cada uma, as vari\'aveis de intera\'c\~ao $y_1$ e $y_2$ (veja o grafo abaixo). No entanto, a natureza da intera\'c\~ao entre $y_1$ e $y_2$ n\~ao \'e modelada. Em particular, n\~ao assumimos uma direcionalidade, i.e., $y_1 \rightarrow y_2$, ou $y_2 \rightarrow y_1$.

\begin{center}
\scalebox{1}{ % x,y
    \begin{tikzpicture}[dgraph]
        
        \node[cloud,opacity=0.5, fill=gray!20, cloud puffs=12, cloud puff arc= 100,
        minimum width=4cm, minimum height=2cm, aspect=1] at (1,-2) {};
        \node[] at (1,-2.5) {\tiny alguma intera\'c\~ao};
        \node[cont] (x1) at (0,-0.5) {$x_1$};
        \node[cont] (x2) at (2,-0.5) {$x_2$};
        \node[cont] (y1) at (0,-2) {$y_1$};
        \node[cont] (y2) at (2,-2) {$y_2$};
        \draw(x1) -- (y1);
        \draw(x2) -- (y2);
        \draw[-,dotted](y1) -- (y2);
\end{tikzpicture}}
\end{center}


\begin{exenumerate}

\item Use as caracteriza\'c\~oes b\'asicas de independ\^encia estat\'istica
\begin{align}
    u \independent v | z &\Longleftrightarrow p(u,v | z) = p(u|z) p(v | z) \label{eq:ind1-alt}\\
    u \independent v | z &\Longleftrightarrow p(u,v | z) = a(u,z) b(v,z) \quad \quad \quad (a(u,z) \ge0, b(v,z) \ge 0) \label{eq:ind2-alt}
\end{align}
para mostrar que $p(y_1,y_2,x_1,x_2)$ satisfaz as seguintes independ\^encias
\begin{align*}
    x_1 \independent x_2 &&  x_1 \independent y_2 \mid y_1, x_2 && x_2 \independent y_1 \mid y_2, x_1
\end{align*}

\begin{solution}
    A fdp/fmp \'e
    $$p(y_1,y_2,x_1,x_2) = p(y_1 | x_1) p(y_2 | x_2) \phi(y_1,y_2) n(x_1,x_2) p(x_1) p(x_2)$$
    
    Para $\mathbf{x_1 \independent x_2}$\\
    Computamos $p(x_1,x_2)$ como
    \begin{align}
        p(x_1,x_2) & = \int p(y_1,y_2,x_1,x_2) \ud y_1 \ud y_2\\
        & =  \int p(y_1 | x_1) p(y_2 | x_2) \phi(y_1,y_2) n(x_1,x_2) p(x_1) p(x_2) \ud y_1 \ud y_2 \\
        & =  n(x_1,x_2) p(x_1) p(x_2) \int p(y_1 | x_1) p(y_2 | x_2) \phi(y_1,y_2) \ud y_1 \ud y_2\\
        & \stackrel{\eqref{eq:ndef-alt}}=  n(x_1,x_2) p(x_1) p(x_2) \frac{1}{n(x_1,x_2)}\\
        & = p(x_1) p(x_2).
    \end{align}
    Como $p(x_1)$ e $p(x_2)$ s\~ao as marginais univariadas de $x_1$ e $x_2$, respectivamente, segue de \eqref{eq:ind1-alt} que $x_1 \independent x_2$.
    
    \vspace{2ex}
    Para $\mathbf{x_1 \independent y_2 \mid y_1, x_2}$\\
    Reescrevemos $p(y_1,y_2,x_1,x_2)$ como
    \begin{align}
        p(y_1,y_2,x_1,x_2) & =  p(y_1 | x_1) p(y_2 | x_2) \phi(y_1,y_2) n(x_1,x_2) p(x_1) p(x_2)\\
        & =  \left[p(y_1 | x_1) p(x_1)  n(x_1,x_2)\right]\left[ p(y_2 | x_2)\phi(y_1,y_2)p(x_2)\right]\\
        & = \phi_A(x_1,y_1,x_2) \phi_B(y_2,y_1,x_2)
    \end{align}
    Com \eqref{eq:ind2-alt}, temos que $x_1 \independent y_2 \mid y_1,x_2$. Note que $p(x_2)$ pode ser associado tanto a $\phi_A$ quanto a $\phi_B$.
    
    \vspace{2ex}
    Para $\mathbf{x_2 \independent y_1 \mid y_2, x_1}$\\ Usamos aqui a mesma abordagem de $x_1 \independent y_2 \mid y_1, x_2$. (Por considera\'c\~oes de simetria, poder\'iamos ver imediatamente que a rela\'c\~ao se sustenta, mas vamos escrev\^e-la por clareza). Reescrevemos $p(y_1,y_2,x_1,x_2)$ como
    \begin{align}
        p(y_1,y_2,x_1,x_2) & =  p(y_1 | x_1) p(y_2 | x_2) \phi(y_1,y_2) n(x_1,x_2) p(x_1) p(x_2)\\
        & =  \left[  p(y_2 | x_2)  n(x_1,x_2) p(x_2)p(x_1) )\right]\left[ p(y_1 | x_1) \phi(y_1,y_2) ]\right)\\
        & = \tilde{\phi}_A(x_2,x_1,y_2) \tilde{\phi}_B(y_1,y_2,x_1)
    \end{align}
    Com \eqref{eq:ind2-alt}, temos que $x_2 \independent y_1 \mid y_2,x_1$. 
    
\end{solution}


\item Existe um mapa perfeito n\~ao direcionado para as independ\^encias satisfeitas por $p(y_1,y_2,x_1,x_2)$?

\begin{solution}
    Escrevemos 
    $$p(y_1,y_2,x_1,x_2) = p(y_1 | x_1) p(y_2 | x_2) \phi(y_1,y_2) n(x_1,x_2) p(x_1) p(x_2)$$
    como uma distribui\'c\~ao de Gibbs
    \begin{align}
        p(y_1,y_2,x_1,x_2) & = \phi_1(y_1,x_1) \phi_2(y_2,x_2) \phi_3(y_1,y_2) \phi_4(x_1,x_2) \quad \quad \text{com} \\
        \phi_1(y_1,x_1) & =p(y_1 | x_1)p(x_1)   \\
        \phi_2(y_2,x_2)& =  p(y_2 | x_2)p(x_2)\\
        \phi_3(y_1,y_2)& =  \phi(y_1,y_2) \\
        \phi_4(x_1,x_2)& = n(x_1,x_2). 
    \end{align}
    Visualiz\'a-lo como um grafo n\~ao direcionado fornece um I-map:
    \begin{center}
        \scalebox{1}{ % x,y
            \begin{tikzpicture}[ugraph]
                \node[cont] (x1) at (0,-0.5) {$x_1$};
                \node[cont] (x2) at (2,-0.5) {$x_2$};
                \node[cont] (y1) at (0,-2) {$y_1$};
                \node[cont] (y2) at (2,-2) {$y_2$};
                \draw(x1) -- (y1);
                \draw(x2) -- (y2);
                \draw(y1) -- (y2);
                \draw(x1) -- (x2);
        \end{tikzpicture}}
    \end{center}
    Enquanto o grafo implica $x_1 \independent y_2 \mid y_1, x_2$ e $x_2 \independent y_1 \mid y_2, x_1$, a independ\^encia $x_1 \independent x_2$ n\~ao \'e representada. Portanto, o grafo n\~ao \'e um mapa perfeito. Note ainda que remover qualquer aresta resultaria em um grafo que n\~ao seria mais um I-map para $\Ind(p)$. Assim, o grafo \'e um I-map m\'inimo para $\Ind(p)$, mas n\~ao podemos obter um I-map perfeito.
    
\end{solution}
\item Existe um mapa perfeito direcionado para as independ\^encias satisfeitas por $p(y_1,y_2,x_1,x_2)$?

\begin{solution} Constru\'imos I-maps direcionados m\'inimos para $p(y_1,y_2,x_1,x_2) = p(y_1, y_2 | x_1, x_2)p(x_1) p(x_2)$ para diferentes ordena\'c\~oes. Veremos que eles n\~ao representam todas as independ\^encias em $\Ind(p)$ e, portanto, n\~ao s\~ao I-maps perfeitos.
    
    Para garantir a independ\^encia n\~ao condicional de $x_1$ e $x_2$, as duas vari\'aveis devem vir primeiro nas ordena\'c\~oes (seja $x_1$ e depois $x_2$ ou vice-versa).
    
    Se usarmos a ordena\'c\~ao $x_1,x_2,y_1,y_2$, e dado que 
    \begin{itemize}
        \item $x_1 \independent x_2$
        \item $y_2 \independent x_1 | y_1, x_2$, que \'e $y_2 \independent \pre(y_2) \setminus \pi | \pi$ para $\pi = (y_1, x_2)$
    \end{itemize}
    est\~ao em $\Ind(p)$, obtemos o seguinte I-map direcionado m\'inimo:
    \begin{center}
        \scalebox{0.9}{ % x,y
            \begin{tikzpicture}[dgraph]
                \node[cont] (x1) at (0,0) {$x_1$};
                \node[cont] (x2) at (2,0) {$x_2$};
                \node[cont] (y1) at (0,-2) {$y_1$};
                \node[cont] (y2) at (2,-2) {$y_2$};
                \draw(x1) -- (y1);
                \draw(x2) -- (y2);
                \draw(x2) -- (y1);
                \draw(y1) -- (y2);
            \end{tikzpicture}
        }
    \end{center}
    O grafo n\~ao representa $x_2 \independent y_1 \mid y_2, x_1$.
    
    Se usarmos a ordena\'c\~ao $x_1,x_2,y_2,y_1$, e dado que
    \begin{itemize}
        \item $x_1 \independent x_2$ 
        \item $y_1 \independent x_2 | x_1, y_2$, que \'e $y_1 \independent \pre(y_1) \setminus \pi | \pi$ para $\pi = (x_1, y_2)$
    \end{itemize}
    est\~ao em $\Ind(p)$, obtemos o seguinte I-map direcionado m\'inimo:
    \begin{center}
        \scalebox{0.9}{ % x,y
            \begin{tikzpicture}[dgraph]
                \node[cont] (x1) at (0,0) {$x_1$};
                \node[cont] (x2) at (2,0) {$x_2$};
                \node[cont] (y1) at (0,-2) {$y_1$};
                \node[cont] (y2) at (2,-2) {$y_2$};
                \draw(x1) -- (y1);
                \draw(x2) -- (y2);
                \draw(y2) -- (y1);
                \draw(x1) -- (y2);
            \end{tikzpicture}
        }
    \end{center}
    O grafo n\~ao representa $x_1 \independent y_2 \mid y_1, x_2$.
    
    Al\'em disso, os grafos implicam uma direcionalidade entre $y_1$ e $y_2$, ou uma influ\^encia direta de $x_1$ em $y_2$, ou de $x_2$ em $y_1$, em contraste com os objetivos originais de modelagem.
    
\end{solution}


\item \emph{(avan\'cado)} O seguinte grafo de fatores representa $p(y_1,y_2,x_1,x_2)$:
\begin{center}
    \scalebox{0.75}{ % x,y
        \begin{tikzpicture}[dgraph]
            \node[] (middle) at (1.5,0) {};
            \node[fact, label=left: $p(x_1)$] (f1) at (0,0) {};
            \node[cont, below=of f1] (x1) {$x_1$};
            \node[fact, label=right: $p(x_2)$] (f2) at (3,0) {};
            \node[cont, below=of f2] (x2) {$x_2$};
            
            \node[fact, below =of x1, label=left: $p(y_1 | x_1)$] (fy1) {};
            \node[fact, below =of x2, label=right: $p(y_2 | x_2)$] (fy2) {};
            
            \node[cont, below =of fy1] (y1) {$y_1$};
            \node[cont, below =of fy2] (y2) {$y_2$};
            
            \node[fact, below= 2 of middle, label=above: $n(x_1 \,x_2)$] (n) {};
            \node[fact, below= 1 of n, label={[label distance=0.25cm]below: $\phi(y_1\,y_2)$}] (phi) {};
            
            \draw[-] (f1) -- (x1);
            \draw[-] (f2) -- (x2);
            \draw  (x1) -- (fy1);
            \draw  (x2) -- (fy2);
            \draw  (x1) -- (n);
            \draw  (x2) -- (n);
            \draw  (fy1) -- (y1);
            \draw  (fy2) -- (y2);
            \draw  (phi) -- (y1);
            \draw  (phi) -- (y2);
            \draw[dashed] (n) -- (y1);
            \draw[dashed] (n) -- (y2);
    \end{tikzpicture}}
\end{center}

Use as regras de separa\'c\~ao para grafos de fatores para verificar que podemos encontrar todas as rela\'c\~oes de independ\^encia. As regras de separa\'c\~ao s\~ao \citep[veja][Se\'c\~ao 4.4.1]{Barber2012}, ou o artigo original de \citet{Frey2003}:\\[1ex]
``Se todos os caminhos estiverem bloqueados, as vari\'aveis s\~ao condicionalmente independentes. Um caminho \'e bloqueado se uma ou mais das seguintes condi\'c\~oes forem satisfeitas:
\begin{enumerate}
    \item Uma das vari\'aveis no caminho est\'a no conjunto de condicionamento.
    \item Uma das vari\'aveis ou fatores no caminho possui duas arestas de entrada que fazem parte do caminho (colisor de vari\'avel ou fator), e nem a vari\'avel ou fator nem qualquer um de seus descendentes est\~ao no conjunto de condicionamento.''
\end{enumerate}

Observa\'c\~oes:
\begin{itemize}
    \item ``uma ou mais das seguintes'' deve ser lido como ``uma das seguintes''. 
    \item ``arestas de entrada'' refere-se a arestas direcionadas de entrada.
    \item os descendentes de um n\'o de vari\'avel ou fator s\~ao todas as vari\'aveis que voc\^e pode alcan\'car seguindo um caminho (contendo arestas direcionadas ou n\~ao direcionadas, mas para arestas direcionadas, todas as dire\'c\~oes devem ser consistentes).
    \item No grafo temos arestas direcionadas tracejadas: elas contam quando voc\^e determina os descendentes, mas n\~ao contribuem para os caminhos. Por exemplo, $y_1$ \'e um descendente do n\'o do fator $n(x_1,x_2)$, mas $x_1 - n - y_2$ n\~ao \'e um caminho.
\end{itemize}

\begin{solution}
    $\mathbf{x_1 \independent x_2}$\\
    Existem dois caminhos de $x_1$ para $x_2$ marcados em vermelho e azul abaixo:
    \begin{center}
        \scalebox{0.75}{ % x,y
            \begin{tikzpicture}[dgraph]
                \node[] (middle) at (1.5,0) {};
                \node[fact, label=left: $p(x_1)$] (f1) at (0,0) {};
                \node[cont, below=of f1] (x1) {$x_1$};
                \node[fact, label=right: $p(x_2)$] (f2) at (3,0) {};
                \node[cont, below=of f2] (x2) {$x_2$};
                
                \node[fact, below =of x1, label=left: $p(y_1 | x_1)$] (fy1) {};
                \node[fact, below =of x2, label=right: $p(y_2 | x_2)$] (fy2) {};
                
                \node[cont, below =of fy1] (y1) {$y_1$};
                \node[cont, below =of fy2] (y2) {$y_2$};
                
                \node[fact, below= 2 of middle, label=above: $n(x_1 \,x_2)$] (n) {};
                \node[fact, below= 1 of n, label={[label distance=0.25cm]below: $\phi(y_1\,y_2)$}] (phi) {};
                
                \draw[-] (f1) -- (x1);
                \draw[-] (f2) -- (x2);
                \draw[blue]  (x1) -- (fy1);
                \draw[blue]  (x2) -- (fy2);
                \draw[red]  (x1) -- (n);
                \draw[red]  (x2) -- (n);
                \draw[blue]  (fy1) -- (y1);
                \draw[blue]  (fy2) -- (y2);
                \draw[blue]  (phi) -- (y1);
                \draw[blue]  (phi) -- (y2);
                \draw[dashed] (n) -- (y1);
                \draw[dashed] (n) -- (y2);
        \end{tikzpicture}}
    \end{center}
    Tanto o caminho azul quanto o vermelho est\~ao bloqueados pela condi\'c\~ao 2. 
    
    $\mathbf{x_1 \independent y_2 \mid y_1, x_2}$\\
    Existem dois caminhos de $x_1$ para $y_2$ marcados em vermelho e azul abaixo:
    \begin{center}
        \scalebox{0.75}{ % x,y
            \begin{tikzpicture}[dgraph]
                \node[] (middle) at (1.5,0) {};
                \node[fact, label=left: $p(x_1)$] (f1) at (0,0) {};
                \node[cont, below=of f1] (x1) {$x_1$};
                \node[fact, label=right: $p(x_2)$] (f2) at (3,0) {};
                \node[contobs, below=of f2] (x2) {$x_2$};
                
                \node[fact, below =of x1, label=left: $p(y_1 | x_1)$] (fy1) {};
                \node[fact, below =of x2, label=right: $p(y_2 | x_2)$] (fy2) {};
                
                \node[contobs, below =of fy1] (y1) {$y_1$};
                \node[cont, below =of fy2] (y2) {$y_2$};
                
                \node[fact, below= 2 of middle, label=above: $n(x_1 \,x_2)$] (n) {};
                \node[fact, below= 1 of n, label={[label distance=0.25cm]below: $\phi(y_1\,y_2)$}] (phi) {};
                
                \draw[-] (f1) -- (x1);
                \draw[-] (f2) -- (x2);
                \draw[blue]  (x1) -- (fy1);
                \draw[red]  (x2) -- (fy2);
                \draw[red]  (x1) -- (n);
                \draw[red]  (x2) -- (n);
                \draw[blue]  (fy1) -- (y1);
                \draw[red]  (fy2) -- (y2);
                \draw[blue]  (phi) -- (y1);
                \draw[blue]  (phi) -- (y2);
                \draw[dashed] (n) -- (y1);
                \draw[dashed] (n) -- (y2);
        \end{tikzpicture}}
    \end{center}
    As vari\'aveis observadas est\~ao marcadas em azul. Para o caminho vermelho, o $x_2$ observado bloqueia o caminho (condi\'c\~ao 1). Note que o n\'o $n(x_1,x_2)$ seria aberto pela condi\'c\~ao 2. O caminho azul tamb\'em est\'a bloqueado pela condi\'c\~ao 1. Em modelos gr\'aficos direcionados, o n\'o $y_1$ estaria aberto, mas aqui, embora a condi\'c\~ao 2 n\~ao se aplique, a condi\'c\~ao 1 ainda se aplica (note o \emph{uma ou mais de...} nas regras de separa\'c\~ao), de modo que o caminho est\'a bloqueado.
    
    $\mathbf{x_2 \independent y_1 \mid y_2, x_1}$\\
    Existem dois caminhos de $x_2$ para $y_1$ marcados em vermelho e azul abaixo:
    \begin{center}
        \scalebox{0.75}{ % x,y
            \begin{tikzpicture}[dgraph]
                \node[] (middle) at (1.5,0) {};
                \node[fact, label=left: $p(x_1)$] (f1) at (0,0) {};
                \node[contobs, below=of f1] (x1) {$x_1$};
                \node[fact, label=right: $p(x_2)$] (f2) at (3,0) {};
                \node[cont, below=of f2] (x2) {$x_2$};
                
                \node[fact, below =of x1, label=left: $p(y_1 | x_1)$] (fy1) {};
                \node[fact, below =of x2, label=right: $p(y_2 | x_2)$] (fy2) {};
                
                \node[cont, below =of fy1] (y1) {$y_1$};
                \node[contobs, below =of fy2] (y2) {$y_2$};
                
                \node[fact, below= 2 of middle, label=above: $n(x_1 \,x_2)$] (n) {};
                \node[fact, below= 1 of n, label={[label distance=0.25cm]below: $\phi(x_1\,x_2)$}] (phi) {};
                
                \draw[-] (f1) -- (x1);
                \draw[-] (f2) -- (x2);
                \draw[blue]  (x1) -- (fy1);
                \draw[red]  (x2) -- (fy2);
                \draw[blue]  (x1) -- (n);
                \draw[blue]  (x2) -- (n);
                \draw[blue]  (fy1) -- (y1);
                \draw[red]  (fy2) -- (y2);
                \draw[red]  (phi) -- (y1);
                \draw[red]  (phi) -- (y2);
                \draw[dashed] (n) -- (y1);
                \draw[dashed] (n) -- (y2);
        \end{tikzpicture}}
    \end{center}
    O mesmo racioc\'inio de antes fornece o resultado.
    
    Finalmente, note que $x_1$ e $x_2$ n\~ao s\~ao independentes dados $y_1$ ou $y_2$ porque o caminho superior atrav\'es de $n(x_1,x_2)$ n\~ao \'e bloqueado sempre que $y_1$ ou $y_2$ s\~ao observados (condi\'c\~ao 2).\\[1ex]
    {\small Cr\'edito: este exemplo \'e discutido no artigo original de B. Frey (Figura 6).}
    
\end{solution}    
\end{exenumerate}